\chapter{Related Work}
\label{chap:related-work}

This chapter discusses existing tools and approaches for Petri net visualization and simulation, and positions the L3S-Offshore-2 Frontend within this landscape.

\section{Desktop Petri Net Tools}
\label{sec:desktop-tools}

\subsection{WoPeD (Workflow Petri Net Designer)}

WoPeD is a widely-used open-source tool for editing and analyzing Workflow Petri nets. It provides:

\begin{itemize}
    \item Graphical editor for place/transition nets
    \item Structural analysis (e.g., place invariants)
    \item Token game simulation
    \item PNML import/export
\end{itemize}

\textbf{Comparison:} Unlike WoPeD, the L3S-Offshore-2 Frontend runs entirely in the browser, enabling collaboration without software installation. However, WoPeD offers more advanced analysis algorithms.

\subsection{CPN Tools}

CPN Tools is a tool for editing, simulating, and analyzing Colored Petri Nets (CPNs). It provides hierarchical modeling and state space analysis.

\textbf{Comparison:} CPN Tools supports more expressive Petri net types but requires desktop installation. The L3S frontend focuses on simpler P/T nets with web-based accessibility.

\section{Web-Based Petri Net Tools}
\label{sec:web-tools}

\subsection{i-love-petrinets (Original Fork)}

The L3S-Offshore-2 Frontend is based on a fork of the \texttt{i-love-petrinets} project, an Angular-based PNML viewer. The original provided:

\begin{itemize}
    \item PNML parsing and basic visualization
    \item Drag-and-drop editing
    \item Place invariant analysis
\end{itemize}

\textbf{Extensions in this project:}
\begin{itemize}
    \item Backend integration for simulation
    \item Animated firing sequence playback
    \item Offshore-specific parameter interface
    \item Production-ready deployment with Docker
    \item Significant UX improvements (zoom, pan, tab state management)
\end{itemize}

\subsection{Cytoscape.js}

Cytoscape.js is a general-purpose graph visualization library that has been extended for Petri nets. However, it lacks native PNML support and Petri net semantics.

\section{Simulation Backends}
\label{sec:simulation-backends}

\subsection{gspn4py}

The L3S backend uses \texttt{gspn4py}, a Python library for Generalized Stochastic Petri Nets. It provides:

\begin{itemize}
    \item Stochastic simulation with exponential delays
    \item Firing sequence generation
    \item Performance analysis
\end{itemize}

\subsection{pm4py}

pm4py is a Python library for process mining that also supports Petri net operations. It was considered during the project's initial phase but the focus shifted to custom simulation for offshore scheduling.

\section{Summary}

The L3S-Offshore-2 Frontend fills a niche for web-based Petri net visualization with domain-specific features for offshore logistics research. While desktop tools offer more advanced analysis, the web-based approach prioritizes accessibility and collaboration.